%-----------------------------------------------------------------------------
% EXPERIMENT DESIGN
%-----------------------------------------------------------------------------
\section{Experiment Design}
\label{sec:experimentdesign}

Before executing the tests in the definitive form described below, we first ran
a lighter version of the same experiments - with fewer repetitions and a smaller
number of sampled parameter configurations - to explore preliminary results and
guide the design of the full test suite. For clarity and conciseness, these
preliminary runs are not reported here, as they constitute a reduced variant of
the tests described in the following subsections.

\subsection{Test 01 - Baseline Model Comparison}
\label{subsec:test01}

\subsection{Test 02 - Per-Agent Sensitivity Analysis (OFAT)}
\label{subsec:test02}

\subsection{Test 03 - \texttt{max\_tokens} Ladder per Agent}
\label{subsec:test03}

\subsection{Test 04 - Compute Expansion (BoN + CoT)}
\label{subsec:test04}

\subsection{Test 05 - Final Pareto Confirmation}
\label{subsec:test05}

%-----------------------------------------------------------------------------

\subsection{Deployment with Docker}
\label{subsec:docker}

\subsubsection{Container Image}

The Sales Agent is packaged as a Docker container to ensure fully reproducible
execution across different server environments. The image is based on
\texttt{python:3.12-slim} and is built in two layers: a dependency layer that
installs the Python packages from \texttt{requirements.txt} (cached unless the
file changes), and an application layer that copies the source code, data files,
and configuration. System-level packages required by Pandas, PyArrow, CodeCarbon,
and Matplotlib (\texttt{build-essential}, \texttt{curl}, \texttt{git}) are
installed via \texttt{apt-get} in the dependency layer.

Two environment variables are set at build time:
\begin{itemize}
    \item \texttt{MPLBACKEND=Agg} — switches Matplotlib to the non-interactive
          Agg backend, enabling chart code execution in headless Linux environments
          without a display server;
    \item \texttt{OLLAMA\_HOST=http://localhost:11434} — sets the default Ollama
          endpoint, which can be overridden at runtime via
          \texttt{-e OLLAMA\_HOST=http://<host-ip>:11434}.
\end{itemize}

The output directory \texttt{/app/runs} is declared as a Docker volume. Benchmark
results are persisted on the host by mounting a local path at container startup:

\begin{verbatim}
docker run --rm --network=host \
    -v /host/path/runs:/app/runs data-agent \
    python evaluation/bulk_runner.py ...
\end{verbatim}

The \texttt{--network=host} flag is used on Linux servers to allow the container
to reach the Ollama process running on the host at \texttt{localhost:11434}.

\subsubsection{Bulk Runner and Ablation Protocol}

The default container entrypoint runs \texttt{bulk\_runner.py}, which implements a
structured ablation study over the agent's hyperparameter space. The protocol
consists of three sequential phases, one per processing step
(\texttt{lookup\_sales\_data}, \texttt{analyzing\_data},
\texttt{create\_visualization}). In each phase, only the target step's
hyperparameters are varied while the other two steps are kept at their default
configuration ($n=1$, fixed temperature). This one-step-at-a-time design isolates
the contribution of each step's hyperparameters to the overall agent performance.

In each phase, $N$ configurations are sampled at random from the search space
defined in \texttt{search\_space.yaml} (with fixed seed 42 for reproducibility),
and the full 20-item benchmark is executed for each configuration. A configurable
think time drawn from an exponential distribution is inserted between consecutive
runs to avoid saturating the LLM API. The full experimental protocol runs
$3 \times 50 = 150$ configurations.

\subsubsection{Hyperparameter Search Space}

The search space for each step is defined in \texttt{search\_space.yaml} and
covers the parameters listed in Table~\ref{tab:searchspace}. Parameters are
specified as fixed scalars, discrete uniform choices, or continuous uniform
intervals. Provider-specific parameters (top-$k$, beam width,
\texttt{no\_repeat\_ngram\_size}) are sampled only for Ollama-based runs and
silently ignored for OpenAI, which does not expose these controls.

\begin{table}[H]
\centering
\arrayrulecolor{black}
\begin{tabular}{|l|l|l|}
\hline
\rowcolor{bluePoli!30}
\textbf{Parameter} & \textbf{Range / choices} & \textbf{Provider} \T\B \\
\hline
\texttt{n} (BoN candidates)         & $\{1, 2, 3, 4\}$                   & All \T\B \\
\hline
\texttt{bon\_param}                  & temperature, top-$p$ (, top-$k$)   & All (top-$k$: Ollama) \T\B \\
\hline
\texttt{cot\_n}                      & $\{1, 2, 3\}$                      & All \T\B \\
\hline
\texttt{temp\_min}                   & $[0.1,\, 0.3]$                     & All \T\B \\
\hline
\texttt{temp\_max}                   & $[0.3,\, 1.0]$                     & All \T\B \\
\hline
\texttt{top\_p\_min}                 & $[0.7,\, 1.0]$                     & All \T\B \\
\hline
\texttt{max\_tokens}                 & $\{1000, 1500, 2000, 3000\}$       & All \T\B \\
\hline
\texttt{top\_k\_min}                 & $\{$null$, 20, 40, 50\}$           & Ollama only \T\B \\
\hline
\texttt{num\_beams}                  & $\{1, 2\}$                         & Ollama only \T\B \\
\hline
\texttt{no\_repeat\_ngram\_size}     & $\{$null$, 2, 3\}$                 & Ollama only \T\B \\
\hline
\end{tabular}
\caption{Hyperparameter search space per step.}
\label{tab:searchspace}
\end{table}

At the end of each phase, \texttt{aggregate\_results.py} consolidates all
per-configuration \texttt{detail\_combined} CSV files into a single summary,
computing mean, minimum, and maximum scores across benchmark items and
configurations. This aggregated output forms the primary dataset for the
experimental analysis presented in Section~\ref{sec:experiments}.
